\documentclass[]{../../../../tex/classes/styledArticle}
\usepackage{../../../../tex/packages/hebrewSupport}


\author{שחר פרץ}
\title{שאלות תרגול בהיסטריה לקראת מבחן הבגרות}
\begin{document}
	\maketitle
	\subsection*{שאלה 8 מתוך בגרות קיץ 2024}
	\textbf{שאלה: }
	
	\begin{enumerate}[A.]
		\item יש הטוענים שההחלטה של בריטניה להעביר את שאלת א''י לאו''ם נבעה מסיקולי מדיניות פנים, ויש הטוענים שהחלטה זו נבעה משיקולי מדיניות חוץ. נמקו את הטענות באמצעות שלוש עובודת היסטוריות (עובדה אחת לפחות לכל אחת מהטענות). 
		\item \begin{itemize}
			\item הסבירו מדוע גם ארה''ב וגם ברה''מ תמכו בהצעת תוכנית החלוקה שהתקבלה באו''ם בכ''ט בנובמבר 1947. 
			\item הצידו סיבה שגרמה לאחת מן המדינות החברות באו''ם להמנע \textit{או} להתנגד על תוכנית זו. 
		\end{itemize}
	\end{enumerate}
	
	\textbf{תשובה: }
	
	\begin{enumerate}[A.]
		\item בשנת 1947 הבריטים החליטו להעביר את השאלה בדבר גודל א''י לידי האו''ם, ארגון האומות המאוחדות, שיחליט מה יהיה עתיד האזור. המעבר התרחש על רקע עימותים, לעיתים מזוינים, בין התנועה הלאומית הציונית שרצתה להקים מדינה בא''י, לבין הערבים בארץ, שלא היו מעוניינים בכך, ובין הבריטים שהתנגדו לחלק מהפעילות היהודית והערבית. בה העת רחשה בעולם ''\textbf{המלחמה הקרה}``, מלחמה ללא לחימה ישירה, בה ארה''ב והגוש המערבי (מדינות עולם ראשון) ברשותאה ניסו להקדים טכנולוגית וגיאופוליטית את הגוש המזרחי וברה''מ בראשו (מדינות עולם שני). 
		
		מלחמת העולם השנייה גררה אחריה בעיה של \textbf{מאות אלפי פליטים יהודים} שנמלטו מהשואה ומצאו עצמם בלי בית. היא גם גררה מחיר כלכלי כבד מבריטניה שהפכה להיות תלויה כלכלית בהארה''ב, ועל־כן חלק מהגוש המערבי. 
		
		\textbf{שאלת א''י} מתייחסת לשאלה מה יעלה בגורל הארץ; האם תוקם מדינה ציונית שתאפשר עליה של יהודים (בינהם פליטים רבים מהשואה), או שמא יתרחש דבר אחר? עד לפני העברת שאלת א''י לאו''ם, בריטניה הגבילה את הפעילות ובפרט הפעילות הציונית בתחומה, ומנעה את מרבית העלייה. 
		
		יש הטוענים שבריטניה העבירה את א''י לאו''ם \textit{בגלל שיקולי מדיניות חוץ}; דוגמה לשיקולים כאלו, לדוגמה, \textbf{השפעת המלחמה הקרה על בריטניה}. היא באה לידי ביטוי, לדוגמה, \textbf{בלחצים שארה''ב והנשיא טרומן הפעיל על ברטניה} לשחרר מארץ ישראל, ומחמת תלותה של בריטניה בארה''ב במהלך המלחהמ הקרה, זו נאלצה לציית. 
		
		מנגד, יש הטוענים שדווקא \textit{מדיניות הפנים היא זו שהכריעה}. עובדה היסטורית אחת התומכת בטענה זו היא \textbf{הלחץ הפנימי שהופעל על הממשלה מבפנים} החברה הבריטית: היו מי מיושבי בריטניה שבגלל השואה ראו את המאבק הציוני כפתרון לבעיית הפליטים היהודים, ועל־כן מאבק צודק, וייתכן ואלו הפיעלו לחץ נכיר דיו על הממשלה הבריטית לשנות את עמדתה בנושא. 
		
		טיעון נוסף לכך שמדיניות הפנים היא זו שהכריעה, הוא \textbf{השיקול הכלכלי}. הטיפול במאבקים הפנימיים בא''י עלה לבריטניה כסף רב, שהיא התקשתה להשיג לאחר מלחמת העולם השנייה שמוטטה את בריטניה כלכלית. 
		
		\textit{שאלות בנוגע לשאלה הזו: }
		\begin{itemize}
			\item אפשר לכתוב כאן על פעילות התנועות הלאומיות? 
			\item זה מפורט דיו, או שב''עובדה`` הכוונה היא לממש אירועים (לדוגמה, אקסודוס ואיך זה השפיע על דעת הקהל בבריטניה)
		\end{itemize}
		
		\item בין השנים 1939-1945 ניטה מלחמת העולם השנייה. במהלכה, \textbf{בריטניה} (הצד המנצח) \textbf{הותשה כלכלית}, נרצחו 6 מיליון יהודים, ומאות אלפי יהודים הפכו \textbf{לפליטים}, מה שבסופה יצר את בעיית הפליטים היהודים – מאות אלפים ללא בית. את אחד הפתרונות לבעיית הפליטים היהודים הציעה \textbf{הציונות} – תנועה לאומית שמטרתה להקים מדינת לאום לעם היהודי בא''י, וראשיתה בסוף המאה ה־19. הפתרון היה עלייה והתיישבות של אותם היהודים בארץ ישראל. בריטניה ששלטה בארץ באותה התקופה, הייתה מעוניינת בשמירה על שטח זה, וכן הערבים שישבו בארץ התנגדו למדינה ציונית. בשנת 1947 החליטה בריטניה להעביר את השאלה בדבר עתיד א''י (הקמת מדינה ציונית, ערבית, בין־לאומית, דו־לאומית, או פתרון אחר) לאו''ם שמום. 
		
		על כן, בשנת 1947 בכ''ט בנובמבר נערכה ההצבעה בנוגע לתמיכה ב''\textbf{תוכנית החלוקה}``, חלוקה של שטח א''י למדינה יהודית ומדינה ערבית, שתי מדינות לשני עמים. לתוכנית החלוקה הייתה תמיכה רבה בקרב היישוב הציוני שראוה כחלום שהתגשם, לעומת התנגדות ערבים שתפסו אותה כתוכנית פסולה המאיימת על שטח שנלקח מהם. בפרט תמכו בהצבעה ארה''ב וברה''מ. 
		
		באותה העת, ארה''ב וברה''מ הובילו שני גושים גיאופוליטיים – הגוש המערבי והגוש הקומוניסטי. הגושים נלחמו במאבק ללא מלחמה ישירה זה בזה לשליטה גיאופוליטית וטכנולוגי במאבק שנקרא ''\textbf{המלחמה קרה}``. 
		
		ברה''מ צפחתה ביישוב היהודי הסוציאליסטי, ובעולים ממזרח אירופה שחלקם הניכר היו קומוניסטיים, \textbf{וציפתה שהמדינה הציונית שתקום תיטה לכיוון הגוש המזרחי}, ובכך תעזור להגדיל את נוכחותה והשפעתה במזרח התיכון, שבאותה העת היה נתון תחת השפעה בריטית וצרפתית בעיקרה (המעצמות הקולוניאליות ששלטו במרבית השטחים באיזור, אם באמצעות קולוניות ואם באמצעות הסכמים עם הקולוניות־לשעבר). על כן, היא הצביעה בעד תוכנית החלוקה. 
		
		ארה''ב הייתה מעוניינת גם היא בהקמת מדינה יהודית ממגוון סיבות, בינהן \textbf{הכוח הרב שהיה לקהילה היהודית בארה''ב}: הן בקלפי והן בכסמים שהזרימו למפלגה הדמוקרטית. הנשיא טרומן עמד לקראת שנת בחירות ורצה בתמיכת הקהילה היהודית, וזו אחת הסיבות שבגינן בחר לתמוך בתוכנית החלוקה. 
		
		נתבונן בבריטניה שנמשנה. הסיבה לכך הייתה שמצד אחד, לא רצתה לפגוע בקשריה עם המדינות הערביות באזור, מה שיקרה אם תתמוך בתוכנית החלוקה, אך היא גם לא יכלה להצביע נגד עמדת ארה''ב בעקבות תלותה הכלכלית. 
	\end{enumerate}
	
	\subsection*{שאלה 11 מתוך בגרות ריץ 2024}
	\textbf{שאלה: }
	
	\begin{enumerate}[A.]
		\item לפניכם שלושה מדפוסי העלייה שנקטה מדינת ישראל בתקופה זו: ''עליית הצלה``, ''עלייה חשאית``, ''עלייה מבוקרת``. 
		\begin{itemize}
			\item בחרו שניים מן הדפוסים, והסבירו שיקול אחד שגרם למדינת ישראל לנקוט כל אחד מן הדפוסים שבחרתם. 
			\item הציגו קושי אחד בביצוע של אחד מדפוסי העלייה האלה. 
		\end{itemize}
	\end{enumerate}
	
%	\textbf{תוכנית: }
%	\begin{enumerate}[A.]
%		\item \begin{itemize}
%			\item הסבר עלייה: סיבות, תקופה, כמויות, רקע של מלחמות ויישוב קטן. 
%			\item הסבר קצת קולוניאליזם + דפוס עליית הצלה
%			\item הדפוס עלייה חשאית
%			\item הדפוס עלייה מבוקרת
%			\item קושי בביצוע עלייה חשאית – זה חשאי
%		\end{itemize}
%		\item \begin{itemize}
%			\item הצגת העלייה וממדיה
%			\item הצגת היהודי החדש וכור ההיתוך
%			\item הצגת התרומה – איחוד ארץ, ייצור, שפה, 
%			\item קושי ראשון, השמדת תרבויות
%			\item קושי שני, ערעור המוסד המשפחתי
%		\end{itemize}
%	\end{enumerate}
	
	\textbf{תשובה: }
	\begin{enumerate}[A.]
		\item 
		בשנת 1948 הוקמה מדינת ישראל, מדינה לאומית לעם היהודי. אחד מהמניעים להקמתה היה התאפשרות עלייה של מאות אלפי פליטים יהודים שחיו באירופה לאחר מלחמת עולם השנייה, בה היטלר השמיד שישה מיליון יהודים. עכב הקמת המדינה, שנעשתה תוך התנגדות של המדינות האסלאמיות במזרח התיכון, התעוררה אנטישמיות מחוזקת במדינות האסלאם, והתרחשו מספר פוגרומים במדינות כגון לוב כנגד היהודים. אותן המדינות היו בתהליך הדה־קולוניזציה בו המדינות השלטות על מדינות האסלאם, הן בעיקר צרפת ובריטניה, יצאו מהן, ונתנו למקומיים לכונן שלטון כרצונן. הפחד משלטון עצמאי שיתמוך במדיניות האנטישמיות, עורר חשש גדול מאוד לגורלם של היהודים באותן הארצות. פרט לכך, היו יהודים במדינות כגון מרוקו ואיראן שבתחילה לא נמצאו בסכנה לחייהם, אף כי חלקם היו מעוניינים בעלייה ממניעים ציוניים – תמיכה ברעיון של מדינת ישראל. היקף העלייה עמד מעל מיליון איש בשנותיה הראשונות של המדינה, שנקלטו ע''י כ־600 אלף תושבים בארץ. 
		
		מדינת ישראל נקטה במדיניות של \textbf{עליות חשאיות}, כאשר השלטון לא הסכים לעזיבה של היהודים מהמדינה, ו/או כאשר לא היו לישראל קשרים עם השלטון במדינה ממנה עולים. כך לדוגמה היה המקרה עם מצריים, שם היה חשש לחייהם של היהודים השוהים במצריים והתעורר צורך שלהם לעלות, על־אף שמצריים הייתה אוייבת של ישראל. המוסד ארגן מבצעים בהם העולים עלו לישראל ועברו את הגבול במחתרת ללא התרת השלטון במדינה ממנו יצאו לפעילות ציונית. 
		
		היקף העלייה הקשה על מדינת ישראל – לא היו מספיק מקומות מגורים, מצב הכלכלה לא היה טוב מספיק בשביל לספק את צרכיהם של כל העולים החדשים, היה מחסור במזון, ורבים מהעולים התקשו להשתלב בשוק העבודה. לא כל עלייה הייתה עליית הצלה – לאחר קום המדינה, רבים מיהודי מרוקו רצו לעלות ארצה, וכן מאות אלפים מהם עשו זאת. שתי עובדות אלו, גרמו לכך שמדינת ישראל פעלה לעיתים במדיניות של \textbf{עלייה סלקטיבית} – כאשר לא היה צורך להציל יהודים ממדינות, הועדפו אנשים צעירים וציוניים שישתלבו במהרה בכלכלה ויתרמו למדינה, ולא מבוגרים או ילדים שיהוו נטל כלכלי והמדינה תצטרך לדאוג להם. כך היה בשנותיה הראשונות של העלייה ממרוקו. 
		
		קושי אחד בביצוע העלייה החשאית היה \textbf{חוסר הרצון של השלטון במדינה ממנה יצאו לשתף פעולה}. הוא התבטא בצורך לבצע מבצעים חשאיים, לעיתים תוך סיכון נפשות, והוא דרש תכנון של המוסד ומשאבים רבים. 
		
		\item 
		בשנת 1948 הוקמה מדינת ישראל – מדינה ציונית לעם היהודי בארץ ישראל. המתיישבים הראשונים בארץ פעלו תחת הרעיון של \textbf{היהודי החדש} – יהודי עובד, פועל אדמה, דובר עברית, ולא איש עסקים הצמוד לשלטון, כמו שהיו מרבית היהודים באירופה. דמות היהודי החדש אומנם הייתה שונה מדמות היהודי האירופאי האשכנזי, אך התבססה עליה שלא במזיד. לאחר קום המדינה, מאות אלפי יהודים עלו מאפריקה וממערב אסיה. תרבותם הייתה שונה מהותית מזו של היהודים שבארץ. 
		
		כדי לשלב את כלל העולים בארץ, הונהגה \textbf{מדיניות כור ההיתוך}. כחלק מהמדיניות, המדינה שאפה להפוך כל עולה ל''יהודי החדש`` – בין אם באמצעות שליחתם לעבודת אדמה וסלילת דרכים, בין עם באמצעות חינוך, ובאמצעים נוספים. מטרתה הייתה לאחד את כלל העולים לכדי עם משותף, עם זהות ציונית משותפת ומאוחדת לכולם. 
		
		מדיניות כור ההיתוך עזרה לעולים, בעיקר לצעירים מבניהם, להשתלב בתוך החברה הישראלית במהרה. היא \textit{תרמה לכלכלה הישראלית} באמצעות העברת עולים למשרות הרלוונטיות, \textit{חזקה את הבטחון} ע''י שליחת הצעירים לצבא, \textit{ואיחדה את העולים} עם המתיישבים. היא \textit{עודדה את למידת השפה העברית}, בהצלחה לא קטנה, מה שהאיץ אף יותר את שילובם של העולים בחברה. בכך, מדיניות כור ההיתוך תרמה לקליטת העלייה. 
		
		עם זאת, מדיניות כור ההיתוך יצרה גם קשיים. קושי אחד היה שדמות היהודי החדש, שנוצרה בראי היהודי האשכנזי, \textbf{לא תאמה כלל את המסורת והתרבות של העולים}. רבים מהם נאלצו לוותר את תרבותם ועל מנהגיהם רבי השנים, וחלק ניכר מהם התקשה לעשות זאת. 
		
		קושי אחר, הוא \textbf{ערעור המוסד המשפחתי}. רבים היו העולים מתרבויות פטריארכליות בהם האב היה ראש המשפחה, והמוסד המשפחתי היה חשוב ביותר. יישום מדיניות כור ההיתוך עודד את השתלבותם המהירה של הצעירים במדינה, בעוד דור ההורים התקשה לקלוט את השפה והתרבות. כך נוצר מצב שבמקום שהאב יהיה זה שמנהיג את הבית ומכלכל אותו, תפקיד זה עובר לבניהם ובנותיהם. באופן דומה, שליחת הילדים לבתי ספר במסגרת חוק חינוך חובה, שם הילדים עברו חינוך ציוני, הייתה קשה בעבור מי שבא מתרבות בה הילד קרוב יותר לבית ולהורים. אזי, מדיניות כור ההיתוך ערערה ושברה את המשפחה כפי שהכירו אותה אותה חלק מהעולים. 
		
	\end{enumerate}
	
	
	
	
	
\end{document}